
En
effet, traitons d'abord le cas \(r=2\). Soient \(g_i(X_0,X_1,X_2)=X_{2-i}^2\) pour \(i=0,1,2\) et \(h_2(X_0,X_1,X_2)=X_1(X_1+1)X_2-X_0\). Alors (iii) est évident et (ii) est 2.7 (i).

Soient \(X\) dans \(P=\mathbb P_k^N\) et \(\omega=\underline{\Ext}^{N-2}_{{O}_P}({O}_X,\Omega^N_{P/k})\). Alors \(h^1(L(-n))=h^1(\omega\otimes L^{-1}(n))\) et \(\omega\) est inversible en dehors de l'ensemble fini \(Z\) des points singuliers de \(X\) (voir FGA 149 \S 6). Soit \(Y\) une section hyperplane normale de \(X\), qui ne contient pas \(Z\). Alors \(\omega_Y(1)\simeq\Omega^1_{Y/k}\) et on a une suite exacte :
\[
0\longrightarrow \omega\otimes L^{-1}(-1)
\longrightarrow \omega\otimes L^{-1}
\longrightarrow \Omega^1_{Y/k}\otimes L_Y^{-1}(-1)
\longrightarrow0.
\]
Or, \(h^1(\Omega^1_{Y/k}\otimes L_Y^{-1}(n))=h^0(L_Y(-n))\) pour tout \(n\), et \(h^0(L_Y(-n))=0\) pour \(n\ge1\) par 2.5. Donc \(\Omega^1_{Y/k}\otimes L_Y^{-1}(-1)\) est 3-régulier et par 1.4, \(\omega\otimes L^{-1}\) est \(\bigl[3+h^1(\omega\otimes L^{-1}(2))\bigr]\)-régulier. Par conséquent, \(h^1(L(-n))=0\) pour \(n\ge2+h^1(L(-2))\)\,; d'où (i).

Supposons maintenant que \(r\ge3\) et raisonnons par récurrence sur \(r\). Soit \(Y\) une section hyperplane normale de \(X\) et considérons la suite exacte :
\[
H^0(L_Y(-n))\longrightarrow H^1(L(-n-1))
\longrightarrow H^1(L(-n))\longrightarrow H^1(L_Y(-n)).
\tag{2.9.1}
\]
Par 1.7, \(\chi(L_Y(n))=\sum_{i=0}^{r-1}a_{i+1}\binom{n+i}{i}\). Donc pour \(n\ge2+h_2(a_{r-2},a_{r-1},a_r)\), l'hypothèse de récurrence donne \(h^1(L_Y(-n))=0\), et comme \(n\ge2\) par (ii) pour \(r=2\), 2.4 donne \(h^0(L_Y(-n))=0\)\,; d'où \(h^1(L(-n-1))=h^1(L(-n))\). Mais, pour \(n\gg0\), \(h^1(L(n))=0\) (voir FGA 149 \S 6)\,; d'où (i).

Comme hypothèse de récurrence, on suppose que les polynômes \(g_i\) et \(h_i\) sont définis pour \(i\le r-1\). Alors pour tout \(n\), \(h^1(L(n))\le h^1(L(n-1))+h_{r-1}(a_{r-2},a_{r-1},a_r)\) grâce à 2.9.1. Donc par (i) 
